\scanpage{73}
$X$ \uline{préschéma}\qquad Conditions équivalentes

\[
\left]\begin{array}{ll}
\text{(i)}&\text{Tt pt fermé}\\
\text{(ii)}&\text{Tt pt admet un voisinage ouvert loc. cpct}\\
\text{(ii bis)}&\rule{6cm}{0.35pt}\ \text{cpct}\\
\text{(ii ter)}&\text{Tout ouvert affine est cpct.}\\
\text{(ii quater)}&\text{(Si }X\text{ quasi-séparé) }\quad X\text{ est loc. cpct}\\
\text{(iii)}&\text{Les anneaux locaux de }X_{\mathrm{réd}}\text{ sont des corps.}
\end{array}\right.
\]

\[
\left\{\begin{array}{ll}
\text{(IV)}&X^{\mathrm{cons}}\longrightarrow X\quad\text{est un \deleted{\unsure{isomor.}} \uline{homéomorphisme} (\unsure{sic|ou} : un isom)}\\
\text{(\unsure{V|VI|IV} bis)}&\text{\uline{Toute partie} loc. constructible \uline{est ouverte}}\\[-1pt]
\deleted{\text{\unsure{(VI bis)}}}&\quad\inserted{\text{ind-constr.}}\\
\text{(V ter)}&\text{(Si }X\text{ qu. séparé) Toute partie \deleted{loc.} constructible est ouverte}\\
\text{(V quater)}&(\hspace{1.8cm})\quad\text{Tout\deleted{e partie} ouvert affine est fermé}
\end{array}\right.
\]

\[
\begin{aligned}
\text{(i)}&\Longleftrightarrow\text{(iii)} &&\text{trivial,}\\
\text{(IV)}&\Longleftrightarrow\text{(V)} &&\text{id}\\
\text{(V)}&\Longleftrightarrow\text{(V bis)} &&\text{id}\\
\text{(V bis)}&\Longleftrightarrow\text{(V ter)}\Longleftrightarrow\text{(V quater)} &&\text{id}\\
\text{(IV)}&\Longrightarrow\text{(\unsure{ii|i})} &&\text{id}\\
&\quad\inserted{\text{et variantes \unsure{bis ter|version bis} etc}}\\
\text{(ii)}&\Longrightarrow\text{(i)} &&\text{id}\\
\text{(i)}&\Longrightarrow\text{(\unsure{V|IV} quater)} &&\text{car l'adhérence est formée}\\
&&&\text{des spécialisations des \unsure{pts} de }U
\end{aligned}
\]
\margin{\uline{NB} Tout espace loc. cpct tot. discontinu muni d'un faisceau d'anneaux comm dont les fibres sont des corps, est un préschéma (loc. cpct réduit \unsure{--|...})}

\uline{Cor}\qquad \uline{Cond. équiv.}\quad (i) $X$ ``loc. compact'' (pas néc. séparé) réduit ;

(ii) Tout Module sur $X$ plat

(ii bis) Tout Module \deleted{\unsure{qu. coh.}} de prés. finie sur un ouvert affine de $X$ est plat

(ii ter) (Si $X$ qu. sép) tt Module qu. coh. sur $X$ \deleted{\unsure{sur ouvert}} est plat.

(iii) $X^{\mathrm{cons}}\longrightarrow X$ est un isom d'espaces annelés.
% archive sheet number: 72
